
\documentclass{article}
\usepackage{colortbl}
\usepackage{makecell}
\usepackage{multirow}
\usepackage{supertabular}

\begin{document}

\newcounter{utterance}

\centering \large Interaction Transcript for game `referencegame', experiment `number\_grids', episode 1 with NEW\_base\_planning\_clp\_chat.
\vspace{24pt}

{ \footnotesize  \setcounter{utterance}{1}
\setlength{\tabcolsep}{0pt}
\begin{supertabular}{c@{$\;$}|p{.15\linewidth}@{}p{.15\linewidth}p{.15\linewidth}p{.15\linewidth}p{.15\linewidth}p{.15\linewidth}}
    \# & $\;$A & \multicolumn{4}{c}{Game Master} & $\;\:$B\\
    \hline

    \theutterance \stepcounter{utterance}  
    & & \multicolumn{4}{p{0.6\linewidth}}{
        \cellcolor[rgb]{0.9,0.9,0.9}{
            \makecell[{{p{\linewidth}}}]{
                \texttt{\tiny{[P1$\langle$GM]}}
                \texttt{You are given three grids, where each of them is 5 by 5 in size.} \\
\texttt{Grids have empty cells marked with "▢" and filled cells marked with "X".} \\
\texttt{Your task is to generate a referring expression that best describes the target grid while distinguishing it from the two other distractor grids.} \\
\texttt{The first grid is the target grid, and the following two grids are the distractors.} \\
\\ 
\texttt{Target grid:} \\
\\ 
\texttt{□ X X X X} \\
\texttt{□ X □ □ X} \\
\texttt{□ X X X X} \\
\texttt{□ X □ □ X} \\
\texttt{□ X X X X} \\
\\ 
\texttt{Distractor grid 1:} \\
\\ 
\texttt{□ X X X X} \\
\texttt{□ X □ □ X} \\
\texttt{□ X X X X} \\
\texttt{□ □ □ □ □} \\
\texttt{□ X X X X} \\
\\ 
\texttt{Distractor grid 2:} \\
\\ 
\texttt{□ X X X X} \\
\texttt{□ X □ □ X} \\
\texttt{□ X □ □ X} \\
\texttt{□ X □ □ X} \\
\texttt{□ X X X X} \\
\\ 
\texttt{Instruction: Describe the target grid.} \\
\texttt{Generate the referring expression starting with the tag "Expression: " for the given target grid. Omit any other text.} \\
            }
        }
    }
    & \\ \\

    \theutterance \stepcounter{utterance}  
    & \multicolumn{4}{p{0.45\linewidth}}{
        \cellcolor[rgb]{0.95,0.95,0.95}{
            \makecell[{{p{\linewidth}}}]{
                \texttt{\tiny{[P1$|$internal]}}
                \texttt{{[}} \\
\texttt{  \{} \\
\texttt{    "role": "user",} \\
\texttt{    "content": "You are given three grids, where each of them is 5 by 5 in size.\textbackslash{}nGrids have empty cells marked with \textbackslash{}"\textbackslash{}u25a2\textbackslash{}" and filled cells marked with \textbackslash{}"X\textbackslash{}".\textbackslash{}nYour task is to generate a referring expression that best describes the target grid while distinguishing it from the two other distractor grids.\textbackslash{}nThe first grid is the target grid, and the following two grids are the distractors.\textbackslash{}n\textbackslash{}nTarget grid:\textbackslash{}n\textbackslash{}n\textbackslash{}u25a1 X X X X\textbackslash{}n\textbackslash{}u25a1 X \textbackslash{}u25a1 \textbackslash{}u25a1 X\textbackslash{}n\textbackslash{}u25a1 X X X X\textbackslash{}n\textbackslash{}u25a1 X \textbackslash{}u25a1 \textbackslash{}u25a1 X\textbackslash{}n\textbackslash{}u25a1 X X X X\textbackslash{}n\textbackslash{}nDistractor grid 1:\textbackslash{}n\textbackslash{}n\textbackslash{}u25a1 X X X X\textbackslash{}n\textbackslash{}u25a1 X \textbackslash{}u25a1 \textbackslash{}u25a1 X\textbackslash{}n\textbackslash{}u25a1 X X X X\textbackslash{}n\textbackslash{}u25a1 \textbackslash{}u25a1 \textbackslash{}u25a1 \textbackslash{}u25a1 \textbackslash{}u25a1\textbackslash{}n\textbackslash{}u25a1 X X X X\textbackslash{}n\textbackslash{}nDistractor grid 2:\textbackslash{}n\textbackslash{}n\textbackslash{}u25a1 X X X X\textbackslash{}n\textbackslash{}u25a1 X \textbackslash{}u25a1 \textbackslash{}u25a1 X\textbackslash{}n\textbackslash{}u25a1 X \textbackslash{}u25a1 \textbackslash{}u25a1 X\textbackslash{}n\textbackslash{}u25a1 X \textbackslash{}u25a1 \textbackslash{}u25a1 X\textbackslash{}n\textbackslash{}u25a1 X X X X\textbackslash{}n\textbackslash{}nInstruction: Describe the target grid.\textbackslash{}nGenerate the referring expression starting with the tag \textbackslash{}"Expression: \textbackslash{}" for the given target grid. Omit any other text."} \\
\texttt{  \},} \\
\texttt{  \{} \\
\texttt{    "role": "user",} \\
\texttt{    "content": "You are a skilled planner. Above you see the game history. What would be the best next action? Openly reason about the next step, then produce your next action exactly as the game expects it preceded by the separator 'ACT:' on its own line. If multiple actions are required, write ACT: first and then all actions."} \\
\texttt{  \}} \\
\texttt{{]}} \\
            }
        }
    }
    & & \\ \\

    \theutterance \stepcounter{utterance}  
    & \multicolumn{4}{p{0.45\linewidth}}{
        \cellcolor[rgb]{0.95,0.95,0.95}{
            \makecell[{{p{\linewidth}}}]{
                \texttt{\tiny{[P1$|$internal]}}
                \texttt{The target grid has a distinct pattern compared to the distractors. Let's analyze the differences:} \\
\\ 
\texttt{{-} In the target grid, the third row is: □ X X X X} \\
\texttt{{-} In distractor 1, the fourth row is: □ □ □ □ □ (all empty)} \\
\texttt{{-} In distractor 2, the third and fourth rows are: □ X □ □ X and □ X □ □ X (both have empty cells in the middle)} \\
\\ 
\texttt{The key distinguishing feature of the target grid is that the third row is fully filled with Xs except for the first cell, and the fourth row has Xs in the second and last positions, with empty cells in the middle.} \\
\\ 
\texttt{But more importantly, the target grid has a full row of Xs in the third row (except first cell), and the fourth row has Xs in the second and last positions, with empty in the middle.} \\
\\ 
\texttt{Wait — let's look again:} \\
\\ 
\texttt{Target:} \\
\texttt{Row 0: □ X X X X} \\
\texttt{Row 1: □ X □ □ X} \\
\texttt{Row 2: □ X X X X} \\
\texttt{Row 3: □ X □ □ X} \\
\texttt{Row 4: □ X X X X} \\
\\ 
\texttt{Distractor 1:} \\
\texttt{Row 0: □ X X X X} \\
\texttt{Row 1: □ X □ □ X} \\
\texttt{Row 2: □ X X X X} \\
            }
        }
    }
    & & \\ \\

    \theutterance \stepcounter{utterance}  
    & \multicolumn{4}{p{0.6\linewidth}}{
        \cellcolor[rgb]{0.9,0.9,0.9}{
            \makecell[{{p{\linewidth}}}]{
                \texttt{\tiny{[P1$\rangle$GM]}}
                \texttt{The target grid has a distinct pattern compared to the distractors. Let's analyze the differences:} \\
\\ 
\texttt{{-} In the target grid, the third row is: □ X X X X} \\
\texttt{{-} In distractor 1, the fourth row is: □ □ □ □ □ (all empty)} \\
\texttt{{-} In distractor 2, the third and fourth rows are: □ X □ □ X and □ X □ □ X (both have empty cells in the middle)} \\
\\ 
\texttt{The key distinguishing feature of the target grid is that the third row is fully filled with Xs except for the first cell, and the fourth row has Xs in the second and last positions, with empty cells in the middle.} \\
\\ 
\texttt{But more importantly, the target grid has a full row of Xs in the third row (except first cell), and the fourth row has Xs in the second and last positions, with empty in the middle.} \\
\\ 
\texttt{Wait — let's look again:} \\
\\ 
\texttt{Target:} \\
\texttt{Row 0: □ X X X X} \\
\texttt{Row 1: □ X □ □ X} \\
\texttt{Row 2: □ X X X X} \\
\texttt{Row 3: □ X □ □ X} \\
\texttt{Row 4: □ X X X X} \\
\\ 
\texttt{Distractor 1:} \\
\texttt{Row 0: □ X X X X} \\
\texttt{Row 1: □ X □ □ X} \\
\texttt{Row 2: □ X X X X} \\
            }
        }
    }
    & & \\ \\

    \theutterance \stepcounter{utterance}  
    & & & \multicolumn{2}{p{0.3\linewidth}}{
        \cellcolor[rgb]{0.9,0.9,0.9}{
            \makecell[{{p{\linewidth}}}]{
                \texttt{\tiny{[GM$|$GM]}}
                \texttt{Invalid generated expression} \\
            }
        }
    }
    & & \\ \\

\end{supertabular}
}

\end{document}
